\documentclass[11pt]{article}

\usepackage[utf8]{inputenc}
\usepackage[T1]{fontenc}
\usepackage{lmodern}
\usepackage{geometry}
\usepackage{amsmath,amssymb,amsfonts,amsthm,mathtools}
\usepackage{bm}
\usepackage{enumitem}
\usepackage{microtype}
\usepackage{hyperref}
\usepackage{titlesec}
\usepackage{booktabs}
\usepackage{array}
\usepackage{setspace}
\usepackage{mathrsfs}
\usepackage{physics}

\geometry{margin=1in}
\hypersetup{colorlinks=true, linkcolor=black, urlcolor=blue, citecolor=black}
\setlist[enumerate]{topsep=0.4em,itemsep=0.35em}
\setlist[itemize]{topsep=0.3em,itemsep=0.25em}
\titleformat{\section}{\large\bfseries}{\thesection.}{0.5em}{}
\titleformat{\subsection}{\normalsize\bfseries}{\thesubsection.}{0.5em}{}
\setstretch{1.06}

\newcommand{\Aether}{\AE{}ther}
\newcommand{\Aflow}{\AE{}ther-flow}
\newcommand{\TheoryName}{\Aflow{} Interpretation of Relativity}
\newcommand{\TheoryProgram}{\Aether{} / \Aflow{} framework}
\newcommand{\Sub}{\mathcal{A}}
\newcommand{\BridgeMetric}{\mathfrak{g}}

\newtheorem{definition}{Definition}
\newtheorem{lemma}{Lemma}
\newtheorem{proposition}{Proposition}
\newtheorem{theorem}{Theorem}
\newtheorem{remark}{Remark}
\newtheorem{corollary}{Corollary}

\title{\textbf{The \TheoryName{}}\\[0.4em]
\large Constraint-Assisted \texorpdfstring{\(U\)}{U}-Sector Quadratic Elimination and Generalized Exact-Closure Conditions}
\author{}
\date{}

\begin{document}

\maketitle

\begin{abstract}
This manuscript carries the constraint-assisted longitudinal \(U\)-sector family through its next fixed derivational task: full flat-background quadratic elimination and extraction of the generalized exact-closure conditions. The scope remains disciplined. The note does not claim a completed substrate derivation of Einsteinian gravity. It computes only the quadratic status of the candidate family already selected in the active research plan.

The new auxiliary scalar \(C\) changes the elimination problem in one precise way. After the coupled \((\chi,C)\) block is integrated out, the scalar source combination \(\alpha X+\beta_I q^I\), with \(X\equiv \partial_0\sigma-\sigma_0\phi\), feeds back into the \(Q\)-sector as a rank-one deformation of the operator \(k^2\delta_{IJ}+\widehat M_{IJ}\). Defining
\[
\Delta_C\equiv 1-\kappa_\theta M_C^2,
\qquad
\mathcal D_C(k^2)
=
\Delta_C
+
\kappa_\theta\,\beta_I
\bigl(k^2\delta+\widehat M\bigr)^{-1\,IJ}
\beta_J,
\]
the reduced scalar bridge channel is governed by three momentum-dependent coefficients \(\mathcal M_C(k^2)\), \(\Gamma_C(k^2)\), and \(\mathcal B_C(k^2)\), and the observer-accessible scalar bridge self-energy becomes
\[
\Pi_\phi^{(C)}(\omega,\mathbf{k})
=
\sigma_0^2
\frac{\mathcal N_C(k^2)}
{\mathcal M_C(k^2)\omega^2+\mathcal B_C(k^2)},
\qquad
\mathcal N_C
=
\mathcal M_C\mathcal B_C-\Gamma_C^2.
\]
Hence the flat-background exact-closure burden is no longer carried by \(\mathcal B(k^2)\) alone. On the generic nonresonant branch, it is carried by the generalized numerator \(\mathcal N_C(k^2)\). Its infrared coefficients \(\mathcal N_{C,0}\) and \(\mathcal N_{C,2}\) give the generalized exact-closure conditions derived here. The separate question of whether those conditions admit a healthy nondegenerate solution is left to the next manuscript.
\end{abstract}

\section{Introduction}

The active exact-closure line of \TheoryName{} is already fixed by the foundations, dynamics, consistency, relativistic-recovery, flow-geometry, and substrate-bridge manuscripts. On the derivational side, the project has now passed through three sharper filters. First, the minimal candidate action was written explicitly. Second, the coefficient-level linearized-consistency note reduced its flat-background exact-closure burden to the scalar operator \(\mathcal B(k^2)\). Third, the minimal stable-sign no-go note and the derivative-mixing branch consistency note showed that neither the minimal branch nor the first nonminimal \(S/Q\) derivative-mixing branch yields the sought nondegenerate flat-background exact-closure solution.

That sequence fixed the next target. The branch-selection note for the constraint-assisted longitudinal \(U\)-sector family showed why the previous no-go decompositions no longer apply once the divergence of the ordered-flow field is sourced by the order and \(Q\) sectors through an auxiliary scalar \(C\). What that note deliberately did not do was the full Phase 3 elimination. The present manuscript performs that elimination.

The result is the exact analogue of the earlier coefficient-level reductions, but with one decisive change. The exact-closure burden no longer collapses directly to the old scalar operator \(\mathcal B(k^2)\) or to the derivative-mixing generalization \(\mathcal B_\rho(k^2)\). Instead it is carried by a new generalized numerator built from the coupled longitudinal \(U\)-sector and \(Q\)-sector response.

\paragraph{Framework and claim boundary.}
\TheoryProgram{} denotes the broader \Aether{} / \Aflow{} framework built on the claims that the \Aether{} is the underlying four-dimensional substrate of reality, the \Aflow{} is its intrinsic ordered motion, observed three-dimensional space is the local experiential slice of that deeper substrate, S-time is the experienced order of change arising from matter, light, and the \Aflow{}, and observed expansion is the three-dimensional appearance of deeper four-dimensional ordered motion. Within that framework, \TheoryName{} denotes the exact-closure theory in which the effective gravitational dynamics are adopted to be exactly Einsteinian with universal matter coupling. In the active sequence, \emph{adoption} means use of that established relativistic dynamics without claiming substrate derivation, while \emph{derivation} is reserved for a first-principles recovery from explicit substrate variables.

The present note stays inside that boundary. It does not promote the constraint-assisted branch to established theory. It determines the branch's full quadratic elimination and the corresponding generalized exact-closure conditions.

\section{Constraint-Assisted Branch and Background Variables}

The constraint-assisted longitudinal \(U\)-sector family modifies the auxiliary sector by
\begin{equation}
\Delta\mathcal L_C
=
C\Bigl(\nabla_AU^A-\alpha\bigl(U^B\partial_BS-\sigma_0\bigr)-\beta_IQ^I\Bigr)
-\frac{M_C^2}{2}C^2.
\label{eq:LC}
\end{equation}
The full branch auxiliary sector is therefore
\begin{equation}
\mathcal L_{\mathrm{aux}}^{(C)}
=
\mathcal L_{\mathrm{aux}}
+\Delta\mathcal L_C,
\label{eq:LauxC}
\end{equation}
where \(\mathcal L_{\mathrm{aux}}\) is the minimal auxiliary sector already defined in the substrate-bridge and coefficient-level notes.

As in the previous flat-background analyses, use the admissible homogeneous bridge background
\begin{equation}
\bar G_{AB}=\delta_{AB},
\qquad
\bar U^A=\delta^A{}_0,
\qquad
\bar S=\sigma_0X^0,
\qquad
\bar Q^I=0,
\label{eq:background}
\end{equation}
for which
\begin{equation}
\bar\BridgeMetric_{AB}=\eta_{AB}=\mathrm{diag}(-1,1,1,1).
\label{eq:etabridge}
\end{equation}
Write perturbations as
\begin{equation}
G_{AB}=\delta_{AB}+H_{AB},
\qquad
U^A=\bar U^A+V^A,
\qquad
S=\sigma_0X^0+\sigma,
\qquad
Q^I=q^I.
\label{eq:perturbations}
\end{equation}
The unit constraint again gives
\begin{equation}
V^0=-\frac{1}{2}H_{00},
\label{eq:V0}
\end{equation}
and the bridge perturbation is
\begin{equation}
h_{00}=-H_{00},
\qquad
h_{0i}=-H_{0i}-2V_i,
\qquad
h_{ij}=H_{ij}.
\label{eq:bridgepert}
\end{equation}
Decompose
\begin{equation}
V_i=V_i^T+\partial_i\chi,
\qquad
\partial_iV_T^i=0,
\label{eq:Vdecomp}
\end{equation}
and define
\begin{equation}
X\equiv \partial_0\sigma-\sigma_0\phi,
\qquad
\Theta\equiv \nabla^2\chi+\frac{1}{2}\partial_0H,
\label{eq:XTheta}
\end{equation}
with \(h_{00}=-2\phi\) and \(H=\delta^{ij}h_{ij}\).

\section{Full Quadratic Auxiliary Lagrangian}

\begin{proposition}[Quadratic auxiliary coefficients of the constraint-assisted branch]
On the flat background \eqref{eq:background}--\eqref{eq:etabridge}, the full quadratic auxiliary Lagrangian is
\begin{align}
\mathcal L_{\mathrm{aux},C}^{(2)}
=\;&
-\frac{m_S^2}{2}X^2
-\frac{\kappa_S}{2}\,\partial_i\sigma\,\partial_i\sigma
-\frac{1}{2}\mu_S^2\sigma^2
-\mu_{SI}\sigma q^I
\nonumber\\
&-\frac{1}{2}q^I\!\left(-\nabla^2\delta_{IJ}+\widehat M_{IJ}\right)\!q^J
-\frac{\kappa_\theta}{2}\Theta^2
+C\bigl(\Theta-\alpha X-\beta_Iq^I\bigr)
-\frac{M_C^2}{2}C^2
\nonumber\\
&-\frac{\kappa_\Omega}{4}\,
\partial_{[i}\!\left(S_{j]}+V^T_{j]}\right)
\partial^{[i}\!\left(S^{j]}+V_T^{j]}\right).
\label{eq:Laux2C}
\end{align}
\end{proposition}

\begin{proof}
The coefficient-level linearized-consistency note already derived the minimal quadratic auxiliary sector, and the branch-selection note already derived the quadratic expansion of the new interaction \eqref{eq:LC}. At this order the only additional term is
\[
\Delta\mathcal L_C^{(2)}
=
C\bigl(\Theta-\alpha X-\beta_Iq^I\bigr)
-\frac{M_C^2}{2}C^2.
\]
Adding that contribution to the previously derived minimal quadratic coefficients yields \eqref{eq:Laux2C}.
\end{proof}

\begin{remark}
The tensor bridge sector remains the ordinary Fierz--Pauli sector of the bridge metric exactly as in the previous notes. All new quadratic structure is carried by the scalar longitudinal \(U\)-sector and its feedback into the \(Q\)-sector.
\end{remark}

\section{Exact Elimination of the Sourced Longitudinal \texorpdfstring{\(U\)}{U}-Sector}

\begin{definition}[Generic nonresonant constraint-assisted branch]
For the present quadratic elimination, the generic nonresonant branch is the part of parameter space satisfying
\begin{equation}
\Delta_C\equiv 1-\kappa_\theta M_C^2\neq 0.
\label{eq:DeltaC}
\end{equation}
The locus \(\Delta_C=0\) is a separate degenerate branch and is not part of the generic elimination derived below.
\end{definition}

\begin{proposition}[On-shell elimination of the \((\chi,C)\) block]
On the decaying flat branch with no harmonic zero mode and \(\Delta_C\neq0\), the coupled \((\chi,C)\) equations imply
\begin{equation}
\Theta
=
\frac{\alpha X+\beta_Iq^I}{\Delta_C},
\qquad
C
=
\frac{\kappa_\theta(\alpha X+\beta_Iq^I)}{\Delta_C}.
\label{eq:ThetaCsolve}
\end{equation}
Substituting back into \eqref{eq:Laux2C} gives the on-shell longitudinal contribution
\begin{equation}
\mathcal L_{\Theta C,\mathrm{on\mbox{-}shell}}^{(2)}
=
-\frac{\kappa_\theta}{2\Delta_C}\,
\bigl(\alpha X+\beta_Iq^I\bigr)^2.
\label{eq:LThetaCon}
\end{equation}
\end{proposition}

\begin{proof}
Varying \eqref{eq:Laux2C} with respect to \(C\) gives
\begin{equation}
\Theta-\alpha X-\beta_Iq^I-M_C^2C=0.
\label{eq:Ceq}
\end{equation}
Varying with respect to \(\chi\) gives
\begin{equation}
\nabla^2\bigl(-\kappa_\theta\Theta+C\bigr)=0.
\label{eq:chieq}
\end{equation}
On the decaying branch, \eqref{eq:chieq} implies \(C=\kappa_\theta\Theta\). Substituting into \eqref{eq:Ceq} yields
\[
\Delta_C\,\Theta=\alpha X+\beta_Iq^I,
\]
which gives \eqref{eq:ThetaCsolve}. Substituting \eqref{eq:ThetaCsolve} back into the \((\Theta,C)\) part of \eqref{eq:Laux2C} gives \eqref{eq:LThetaCon}.
\end{proof}

\begin{corollary}[Unchanged transverse \(U\)-sector elimination]
On the same decaying branch, the transverse equation is unchanged:
\begin{equation}
V_T^i=-S^i.
\label{eq:VTsolve}
\end{equation}
Consequently the \(\kappa_\Omega\) sector contributes no residual flat-background bridge self-energy at quadratic order.
\end{corollary}

\begin{proof}
The new \(C\)-interaction contains no \(V_T^i\). Therefore the transverse \(U\)-sector variation is exactly the same as in the minimal and derivative-mixing branches.
\end{proof}

After eliminating \((\chi,C)\), the scalar/\(Q\) part of \eqref{eq:Laux2C} becomes
\begin{align}
\mathcal L_{\sigma q,C}^{(2)}
=\;&
-\frac{1}{2}\left(m_S^2+\frac{\alpha^2\kappa_\theta}{\Delta_C}\right)X^2
-\frac{1}{2}\sigma\Bigl(\kappa_S(-\nabla^2)+\mu_S^2\Bigr)\sigma
-\mu_{SI}\sigma q^I
\nonumber\\
&-\frac{1}{2}q^I
\left[
\mathcal A_{IJ}(-\nabla^2)
+\frac{\kappa_\theta}{\Delta_C}\beta_I\beta_J
\right]
q^J
-\frac{\alpha\kappa_\theta}{\Delta_C}\,X\,\beta_Iq^I,
\label{eq:LsigmaqC}
\end{align}
where
\begin{equation}
\mathcal A_{IJ}(-\nabla^2)
\equiv
(-\nabla^2)\delta_{IJ}+\widehat M_{IJ}.
\label{eq:Aoperator}
\end{equation}

\section{Elimination of the \texorpdfstring{\(Q\)}{Q}-Sector and Reduced Scalar Action}

\begin{definition}[Constraint-assisted response operators]
Define the scalar operators
\begin{align}
\mathcal G(-\nabla^2)
&\equiv
\beta_I\mathcal A^{-1\,IJ}(-\nabla^2)\beta_J,
\label{eq:Goperator}
\\
\mathcal Z(-\nabla^2)
&\equiv
\mu_{SI}\mathcal A^{-1\,IJ}(-\nabla^2)\beta_J,
\label{eq:Zoperator}
\\
\Sigma_q(-\nabla^2)
&\equiv
\mu_{SI}\mathcal A^{-1\,IJ}(-\nabla^2)\mu_{SJ},
\label{eq:Sigmaoperator}
\\
\mathcal B(-\nabla^2)
&\equiv
\kappa_S(-\nabla^2)+\mu_S^2-\Sigma_q(-\nabla^2),
\label{eq:Boperator}
\\
\mathcal D_C(-\nabla^2)
&\equiv
\Delta_C+\kappa_\theta\mathcal G(-\nabla^2).
\label{eq:DCoperator}
\end{align}
\end{definition}

\begin{lemma}[Rank-one inverse of the deformed \(Q\)-operator]
Let
\begin{equation}
\mathcal A^{(C)}_{IJ}(-\nabla^2)
\equiv
\mathcal A_{IJ}(-\nabla^2)
+\frac{\kappa_\theta}{\Delta_C}\beta_I\beta_J.
\label{eq:ACoperator}
\end{equation}
On the generic nonresonant branch, its inverse is
\begin{equation}
\bigl(\mathcal A^{(C)}\bigr)^{-1\,IJ}
=
\mathcal A^{-1\,IJ}
-\frac{\kappa_\theta}{\mathcal D_C}\,
\mathcal A^{-1\,IK}\beta_K\beta_L\mathcal A^{-1\,LJ}.
\label{eq:ACinverse}
\end{equation}
\end{lemma}

\begin{proof}
Equation \eqref{eq:ACoperator} is a rank-one deformation of \(\mathcal A_{IJ}\). The Sherman--Morrison identity gives
\[
\bigl(\mathcal A^{(C)}\bigr)^{-1}
=
\mathcal A^{-1}
-\frac{\frac{\kappa_\theta}{\Delta_C}\,\mathcal A^{-1}\beta\beta^\top\mathcal A^{-1}}
{1+\frac{\kappa_\theta}{\Delta_C}\,\beta^\top\mathcal A^{-1}\beta}.
\]
Using \(\beta^\top\mathcal A^{-1}\beta=\mathcal G\) and \(\mathcal D_C=\Delta_C+\kappa_\theta\mathcal G\) yields \eqref{eq:ACinverse}.
\end{proof}

\begin{proposition}[Exact \(Q\)-sector elimination]
The \(q^I\) Euler--Lagrange equation derived from \eqref{eq:LsigmaqC} is
\begin{equation}
\mathcal A_{IJ}^{(C)}(-\nabla^2)\,q^J
=
-\mu_{SI}\sigma
-\frac{\alpha\kappa_\theta}{\Delta_C}\,\beta_I X.
\label{eq:qeqC}
\end{equation}
Hence
\begin{equation}
q^I
=
-\mathcal A^{-1\,IJ}\mu_{SJ}\sigma
-\frac{\kappa_\theta}{\mathcal D_C}\,
\mathcal A^{-1\,IJ}\beta_J\bigl(\alpha X-\mathcal Z\,\sigma\bigr).
\label{eq:qsolveC}
\end{equation}
\end{proposition}

\begin{proof}
Equation \eqref{eq:qeqC} is immediate from \eqref{eq:LsigmaqC}. Applying \eqref{eq:ACinverse} to the source on the right-hand side gives \eqref{eq:qsolveC}.
\end{proof}

\begin{proposition}[Reduced scalar bridge channel]
After eliminating \(q^I\), the scalar bridge sector becomes
\begin{equation}
\mathcal L_{\sigma,\mathrm{eff},C}^{(2)}
=
-\frac{1}{2}\,
X\,\mathcal M_C(-\nabla^2)\,X
+\sigma\,\Gamma_C(-\nabla^2)\,X
-\frac{1}{2}\,
\sigma\,\mathcal B_C(-\nabla^2)\,\sigma,
\label{eq:LsigmaeffC}
\end{equation}
with
\begin{align}
\mathcal M_C(-\nabla^2)
&=
m_S^2+\frac{\alpha^2\kappa_\theta}{\mathcal D_C(-\nabla^2)},
\label{eq:Moperator}
\\
\Gamma_C(-\nabla^2)
&=
\frac{\alpha\kappa_\theta\,\mathcal Z(-\nabla^2)}
{\mathcal D_C(-\nabla^2)},
\label{eq:Gammaoperator}
\\
\mathcal B_C(-\nabla^2)
&=
\mathcal B(-\nabla^2)
+\frac{\kappa_\theta\,\mathcal Z(-\nabla^2)^2}
{\mathcal D_C(-\nabla^2)}.
\label{eq:BCoperator}
\end{align}
\end{proposition}

\begin{proof}
Substituting \eqref{eq:qsolveC} into \eqref{eq:LsigmaqC} and using \eqref{eq:ACinverse} gives
\[
\frac{1}{2}
\bigl(\mu_{SI}\sigma+\tfrac{\alpha\kappa_\theta}{\Delta_C}\beta_I X\bigr)
\bigl(\mathcal A^{(C)}\bigr)^{-1\,IJ}
\bigl(\mu_{SJ}\sigma+\tfrac{\alpha\kappa_\theta}{\Delta_C}\beta_J X\bigr)
\]
for the on-shell \(Q\)-sector contribution. Evaluating the three resulting terms with \eqref{eq:ACinverse} yields the coefficients \eqref{eq:Moperator}--\eqref{eq:BCoperator}, which gives \eqref{eq:LsigmaeffC}.
\end{proof}

\begin{corollary}[Minimal-branch reversion]
If \(\alpha=\beta_I=0\), then
\begin{equation}
\mathcal M_C=m_S^2,
\qquad
\Gamma_C=0,
\qquad
\mathcal B_C=\mathcal B,
\label{eq:reversioncheck}
\end{equation}
so the reduced scalar bridge channel contracts back to the minimal coefficient-level result.
\end{corollary}

\section{Generalized Bridge Self-Energy and Exact-Closure Conditions}

\begin{proposition}[Generalized scalar bridge self-energy]
Up to a boundary term, the mixed term in \eqref{eq:LsigmaeffC} satisfies
\begin{equation}
\sigma\,\Gamma_C(-\nabla^2)\,X
\simeq
-\sigma_0\,\sigma\,\Gamma_C(-\nabla^2)\,\phi.
\label{eq:boundarydrop}
\end{equation}
After integrating out \(\sigma\), the observer-accessible scalar bridge correction is therefore
\begin{equation}
\Delta\mathcal L_{\phi,C}^{(2)}(\omega,\mathbf{k})
=
-\frac{1}{2}\,
\Pi_\phi^{(C)}(\omega,\mathbf{k})\,|\phi(\omega,\mathbf{k})|^2,
\label{eq:DeltaphiC}
\end{equation}
with
\begin{equation}
\Pi_\phi^{(C)}(\omega,\mathbf{k})
=
\sigma_0^2\,
\frac{\mathcal N_C(k^2)}
{\mathcal M_C(k^2)\omega^2+\mathcal B_C(k^2)},
\label{eq:PiphiC}
\end{equation}
where
\begin{equation}
\mathcal N_C(k^2)
\equiv
\mathcal M_C(k^2)\mathcal B_C(k^2)-\Gamma_C(k^2)^2.
\label{eq:NCdef}
\end{equation}
Equivalently,
\begin{equation}
\mathcal N_C(k^2)
=
m_S^2\mathcal B(k^2)
+\frac{\kappa_\theta}
{\mathcal D_C(k^2)}
\Bigl(\alpha^2\mathcal B(k^2)+m_S^2\mathcal Z(k^2)^2\Bigr).
\label{eq:NCcompact}
\end{equation}
\end{proposition}

\begin{proof}
Because \(\Gamma_C(-\nabla^2)\) depends only on the spatial operator, it commutes with \(\partial_0\). Thus
\[
\sigma\,\Gamma_C\,\partial_0\sigma
=
\frac{1}{2}\partial_0\!\left(\sigma\,\Gamma_C\,\sigma\right),
\]
which is a boundary term, giving \eqref{eq:boundarydrop}. The remaining quadratic \(\sigma\)-action then has the same structure as in the earlier coefficient-level notes, but with the replacement
\[
m_S^2\mapsto \mathcal M_C,
\qquad
\mu_{SI}\mapsto \Gamma_C.
\]
Integrating out \(\sigma\) therefore yields \eqref{eq:PiphiC}. Expanding \(\mathcal N_C=\mathcal M_C\mathcal B_C-\Gamma_C^2\) with \eqref{eq:Moperator}--\eqref{eq:BCoperator} gives \eqref{eq:NCcompact}.
\end{proof}

\begin{remark}
The exact-closure burden is now carried by the generalized numerator \(\mathcal N_C(k^2)\), not by \(\mathcal B(k^2)\) alone. This is the precise sense in which the constraint-assisted branch changes the Phase 3 problem.
\end{remark}

To extract the low-momentum conditions, expand
\begin{align}
\mathcal G(k^2)
&=
G_0-G_2k^2+\mathcal O(k^4),
\label{eq:Gexpand}
\\
\mathcal Z(k^2)
&=
Z_0-Z_2k^2+\mathcal O(k^4),
\label{eq:Zexpand}
\\
\mathcal B(k^2)
&=
\mathcal B_0+\mathcal B_2k^2+\mathcal O(k^4),
\label{eq:Bexpand}
\end{align}
with
\begin{align}
G_0
&=
\beta_I\widehat M^{-1\,IJ}\beta_J,
\qquad
G_2
=
\beta_I\widehat M^{-2\,IJ}\beta_J,
\label{eq:Gcoeffs}
\\
Z_0
&=
\mu_{SI}\widehat M^{-1\,IJ}\beta_J,
\qquad
Z_2
=
\mu_{SI}\widehat M^{-2\,IJ}\beta_J,
\label{eq:Zcoeffs}
\\
\mathcal B_0
&=
\mu_S^2-\mu_{SI}\widehat M^{-1\,IJ}\mu_{SJ},
\label{eq:B0}
\\
\mathcal B_2
&=
\kappa_S+\mu_{SI}\widehat M^{-2\,IJ}\mu_{SJ}.
\label{eq:B2}
\end{align}
Define also
\begin{equation}
\mathcal D_{C,0}
\equiv
\Delta_C+\kappa_\theta G_0
=
1-\kappa_\theta M_C^2+\kappa_\theta\beta_I\widehat M^{-1\,IJ}\beta_J.
\label{eq:DC0}
\end{equation}

\begin{proposition}[Infrared coefficients of the generalized numerator]
On the branch with \(\mathcal D_{C,0}\neq0\), the generalized numerator expands as
\begin{equation}
\mathcal N_C(k^2)
=
\mathcal N_{C,0}
+\mathcal N_{C,2}k^2
+\mathcal O(k^4),
\label{eq:NCexpand}
\end{equation}
where
\begin{align}
\mathcal N_{C,0}
&=
m_S^2\mathcal B_0
+\frac{\kappa_\theta}{\mathcal D_{C,0}}
\bigl(\alpha^2\mathcal B_0+m_S^2Z_0^2\bigr),
\label{eq:NC0}
\\
\mathcal N_{C,2}
&=
m_S^2\mathcal B_2
+\frac{\kappa_\theta}{\mathcal D_{C,0}}
\bigl(\alpha^2\mathcal B_2-2m_S^2Z_0Z_2\bigr)
+\frac{\kappa_\theta^2G_2}{\mathcal D_{C,0}^2}
\bigl(\alpha^2\mathcal B_0+m_S^2Z_0^2\bigr).
\label{eq:NC2}
\end{align}
\end{proposition}

\begin{proof}
Using \eqref{eq:Gexpand} gives
\[
\mathcal D_C(k^2)
=
\mathcal D_{C,0}-\kappa_\theta G_2k^2+\mathcal O(k^4),
\]
and therefore
\[
\frac{1}{\mathcal D_C(k^2)}
=
\frac{1}{\mathcal D_{C,0}}
+\frac{\kappa_\theta G_2}{\mathcal D_{C,0}^2}k^2
+\mathcal O(k^4).
\]
Substituting this together with \eqref{eq:Gexpand}--\eqref{eq:Bexpand} into \eqref{eq:NCcompact} gives \eqref{eq:NC0} and \eqref{eq:NC2}.
\end{proof}

\begin{theorem}[Generalized flat-background exact-closure conditions]
On the generic nonresonant constraint-assisted branch with \(\mathcal D_{C,0}\neq0\), quadratic infrared exact closure requires
\begin{equation}
\mathcal N_{C,0}=0,
\qquad
\mathcal N_{C,2}=0.
\label{eq:NCconditions}
\end{equation}
Equivalently, the generalized coefficient conditions are
\begin{equation}
\bigl(m_S^2\mathcal D_{C,0}+\alpha^2\kappa_\theta\bigr)\mathcal B_0
+m_S^2\kappa_\theta Z_0^2
=
0,
\label{eq:gencond0}
\end{equation}
and
\begin{equation}
\bigl(m_S^2\mathcal D_{C,0}+\alpha^2\kappa_\theta\bigr)\mathcal B_2
-2m_S^2\kappa_\theta Z_0Z_2
-m_S^2\kappa_\theta G_2\mathcal B_0
=
0.
\label{eq:gencond2}
\end{equation}
\end{theorem}

\begin{proof}
If \(\mathcal N_{C,0}\neq0\), then \eqref{eq:PiphiC} contains an unsuppressed \(p^0\) scalar bridge correction. If \(\mathcal N_{C,0}=0\) but \(\mathcal N_{C,2}\neq0\), then an unsuppressed \(k^2\) non-Einsteinian scalar correction remains. Therefore quadratic infrared exact closure requires \eqref{eq:NCconditions}. Multiplying \eqref{eq:NC0}=0 by \(\mathcal D_{C,0}\) gives \eqref{eq:gencond0}. Then substituting \eqref{eq:gencond0} into \eqref{eq:NC2}=0 yields \eqref{eq:gencond2}.
\end{proof}

\begin{remark}
Equations \eqref{eq:gencond0}--\eqref{eq:gencond2} complete the fixed task for the constraint-assisted branch: they are the generalized exact-closure conditions replacing the old \(\mathcal B_0=0\), \(\mathcal B_2=0\) pair. They do \emph{not} yet decide whether a healthy nondegenerate solution exists. That requires the next-stage sign and pole analysis of \(\mathcal D_C\), \(\mathcal M_C\), and \(\mathcal B_C\).
\end{remark}

\section{Conclusion}

The constraint-assisted longitudinal \(U\)-sector family has now been carried through its full flat-background quadratic elimination. The \((\chi,C)\) sector does not collapse to a vanishing bracket as it did in the previous branches. Instead it feeds back into the scalar bridge channel as a rank-one deformation of the \(Q\)-sector operator.

That change is now fully explicit. After elimination of the coupled auxiliary system, the scalar bridge self-energy is governed by the generalized numerator \(\mathcal N_C(k^2)\), not by the old operator \(\mathcal B(k^2)\) alone. On the generic nonresonant branch, the quadratic exact-closure conditions are the two infrared coefficient relations \eqref{eq:gencond0}--\eqref{eq:gencond2}.

This completes the immediate derivational task that the branch-selection note left open. The next task is sharper and narrower: determine whether the new conditions admit any healthy nondegenerate branch or whether the sourced longitudinal \(U\)-sector again contracts onto a resonant or otherwise degenerate solution. Until that is settled, exact closure remains the benchmark and reversion theory of the project.

\end{document}
